\scanpage{5}
\mspage{4}
\unsure{pour} \deleted{\unsure{le}}\unsure{tt|\ensuremath{\forall}} $A\in\unsure{\mathscr{L}|\mathcal{L}}$, on a $A^\circ\ne\varnothing$. (Donc l'ens d'indices de la partition de $3^\circ)\,\gamma$ est $\mathscr{L}$ lui-\ensuremath{=})

\deleted{\uline{Proposition 2} Soient $\mathscr{L},\Psi$ deux familles préadmissibles. Conditions équivalentes :}

% slanted
\margin{Souvent il est plus commode de dire que c'est un \uline{préordre} sur $S$ -- $I$ est \unsure{l'ens.} ordonné associé, $\unsure{R|\mathscr{R}}_{\mathscr{L}}$ la \unsure{rel.} \unsure{d'éq.} \unsure{d'équivalence|dite} \unsure{associée}, \unsure{pour|par} le \unsure{préordre|préordonné}. Les $A\in\mathscr{L}$ sont les parties fermées \uline{irréd.}}
% slanted
\margin{\uline{NB} Se donner une famille \uline{pré}admissible, c'est se donner un système comme dessus (\deleted{\unsure{qui corres}} \inserted{où} $I=\mathscr{L}'=\{A\in\mathscr{L}\mid A^\circ\ne\varnothing\}$) et de plus, un ens. \unsure{fini|fin.} \deleted{de \unsure{parti}} $\Psi$ de parties fermées de $I$ (\unsure{correspondant} aux éléments de $\mathscr{L}\setminus\mathscr{L}'$)}

Se donner dans $L$ une famille admissible \inserted{$\mathscr{L}$} de parties, revient au $=$ que de se donner

a) Un ens ordonné $I$ \qquad (une copie conforme de $\mathscr{L}$)

b) Une partie $S$ de $L$ \qquad ($S=|\mathscr{L}|$)

c) Une application surjective $\unsure{\varphi|\phi}:S\unsure{\text{---}|\longrightarrow}I$ \quad ($x\mapsto A(x)$)

Prenant sur $I$ la top. associée à l'ordre (où les fermés sont les réunions quel\uline{conques} d'ens. $I_{\leq i}$, les ouverts les réunions quelc. d'ens $I_{\geq i}$) -- donc $\uline{\ensuremath{J\subset I}}$ est fermé ssi $\deleted{\unsure{\forall}}i\in J,\ j\leq i\Rightarrow j\in J$) et sur $S$ la top. image inverse (qui est aussi celle définie par le préordre induit de l'ordre de $I$ via $\varphi$), \deleted{les} on récupère $\deleted{\unsure{P}}\mathscr{L}$ comme
\[
 \mathscr{L}=\left\{\overline{\varphi^{-1}(i)}=\varphi^{-1}
 \bigl(\overline{\{i\}}=I_{\leq i}\bigr)\right\}_{i\in I}
\]
(indexation fidèle). Les parties fermées de $S$ sont les réunions quelconques \deleted{de} ``strates fermées'' i.e d'éléments de $A$, elles correspondent \deleted{\unsure{d|1} à \unsure{un|un*}}\inserted{biuniv.} aux parties fermées \inserted{$J$} de $I$ ($\unsure{\simeq|\cong}$ \deleted{de} $\mathscr{L}$), \inserted{par transp. inv. par $\varphi$} \deleted{les parties ouvertes sont les correspondants aux}
\[
\deleted{\unsure{X}}A_J=\bigcup_{i\unsure{\leq|\in}J}A_i
 =\underset{\substack{\text{si on identifie}\\I=\mathscr{L},\ \text{donc }J\subset\mathscr{L}}}{\bigcup_{A\in J} A}
 =\underbrace{\deleted{\unsure{\bigcup}}\bigsqcup_{i\in J}A_i^\circ}_{\substack{\text{réunion}\\\text{disjointe}\\\text{de parties}\\\text{non vides}}}
\]
\uuline{NB} \deleted{Si on regarde les}

De $=$, les parties ouvertes sont celles de
